\documentclass{article}
\usepackage{amsmath, amssymb, amsthm}

\newtheorem{theorem}{Theorem}
\newtheorem{proposition}[theorem]{Proposition}
\newtheorem{lemma}[theorem]{Lemma}
\newtheorem{corollary}[theorem]{Corollary}
\theoremstyle{definition}
\newtheorem{assumption}{Assumption}
\theoremstyle{remark}
\newtheorem{remark}{Remark}

\begin{document}

\section*{Discrete-Time Convergence of KL-Regularized\\ Projected GDA for Gaussian Policies}

\subsection*{1. Setup and Notation}

Player 1 and Player 2 select Gaussian policies $\pi_1 = \mathcal{N}(\mu_x, \sigma_x^2)$ and
$\pi_2 = \mathcal{N}(\mu_y, \sigma_y^2)$. Collect the parameters into
$\boldsymbol{\theta} = (\mu_x, \sigma_x, \mu_y, \sigma_y)^{T}$ and restrict them to
\begin{equation}
\Theta = \left\{ (\mu_x, \sigma_x, \mu_y, \sigma_y) \in \mathbb{R}^4 \;\middle|\;
\sigma_x \ge \sigma_{\min}, \; \sigma_y \ge \sigma_{\min} \right\},
\qquad \sigma_{\min} > 0 \text{ fixed.}
\end{equation}
$\Theta$ is nonempty, closed and convex (an intersection of two half-spaces).

The \emph{magnet} (anchor) policy is a fixed
$\bar{\boldsymbol{\theta}} = (\bar\mu_x, \bar\sigma_x, \bar\mu_y, \bar\sigma_y)^{T} \in \Theta$,
inducing $\bar\pi_1 = \mathcal{N}(\bar\mu_x, \bar\sigma_x^2)$ and
$\bar\pi_2 = \mathcal{N}(\bar\mu_y, \bar\sigma_y^2)$. Throughout, $\bar\sigma_x, \bar\sigma_y \ge \sigma_{\min}$.

\paragraph{Notation.} Iterates are indexed by $k$ ($\boldsymbol{\theta}_k$); the magnet carries a bar and
never an index. The strong-monotonicity modulus of the regularizer is written $\rho_R$ (not $\mu_{\mathrm{KL}}$)
so that the symbol $\mu$ is reserved for policy means.

\paragraph{Operators.} With
$U(\boldsymbol{\theta}) = \mathbb{E}_{X \sim \pi_1, Y \sim \pi_2}[f(X,Y)]$ the expected payoff
(Player 1 minimizing, Player 2 maximizing), define the game operator and the regularizer operator
\begin{align}
V(\boldsymbol{\theta}) &= \begin{pmatrix}
\nabla_{\mu_x} U(\boldsymbol{\theta}) \\ \nabla_{\sigma_x} U(\boldsymbol{\theta}) \\
-\nabla_{\mu_y} U(\boldsymbol{\theta}) \\ -\nabla_{\sigma_y} U(\boldsymbol{\theta})
\end{pmatrix},
&
R(\boldsymbol{\theta}) &= \begin{pmatrix}
\nabla_{\mu_x} D_{\mathrm{KL}}(\pi_1 \parallel \bar\pi_1) \\
\nabla_{\sigma_x} D_{\mathrm{KL}}(\pi_1 \parallel \bar\pi_1) \\
\nabla_{\mu_y} D_{\mathrm{KL}}(\pi_2 \parallel \bar\pi_2) \\
\nabla_{\sigma_y} D_{\mathrm{KL}}(\pi_2 \parallel \bar\pi_2)
\end{pmatrix},
\end{align}
and the regularized operator $F(\boldsymbol{\theta}) = V(\boldsymbol{\theta}) + \lambda R(\boldsymbol{\theta})$
with penalty weight $\lambda > 0$.

The sign pattern of $R$ is the crux of the argument and deserves emphasis: Player 1 minimizes
$U + \lambda D_{\mathrm{KL}}(\pi_1 \parallel \bar\pi_1)$ while Player 2 maximizes
$U - \lambda D_{\mathrm{KL}}(\pi_2 \parallel \bar\pi_2)$, so Player 2's block of the
descent-ascent operator is
\begin{equation}
-\nabla_{(\mu_y,\sigma_y)}\!\left[U - \lambda D_{\mathrm{KL}}(\pi_2 \parallel \bar\pi_2)\right]
= -\nabla_{(\mu_y,\sigma_y)} U + \lambda \nabla_{(\mu_y,\sigma_y)} D_{\mathrm{KL}}(\pi_2 \parallel \bar\pi_2).
\end{equation}
The min-max sign flip cancels against the sign of the penalty, and \emph{both} blocks of $\lambda R$ enter with a
$+$ sign. Hence $R = \nabla \Phi$ is the gradient of the separable convex function
\begin{equation}
\Phi(\boldsymbol{\theta}) = D_{\mathrm{KL}}(\pi_1 \parallel \bar\pi_1)
+ D_{\mathrm{KL}}(\pi_2 \parallel \bar\pi_2),
\end{equation}
and is therefore monotone, whereas $V$ need not be.

Updates follow projected gradient descent-ascent (PGDA) with step size $\eta > 0$:
\begin{equation}\label{eq:pgda}
\boldsymbol{\theta}_{k+1} = \Pi_\Theta \big( \boldsymbol{\theta}_k - \eta F(\boldsymbol{\theta}_k) \big),
\end{equation}
where $\Pi_\Theta$ is the Euclidean projection onto $\Theta$.

\subsection*{2. Assumptions}

Only two things are assumed; every property of the KL regularizer is computed in Section 3.

\begin{assumption}[Regularity of the game operator]\label{as:lip}
$V$ is $L_V$-Lipschitz continuous on $\Theta$:
$\|V(\boldsymbol{\theta}) - V(\boldsymbol{\theta}')\| \le L_V \|\boldsymbol{\theta} - \boldsymbol{\theta}'\|$
for all $\boldsymbol{\theta}, \boldsymbol{\theta}' \in \Theta$.
\end{assumption}

Note that no monotonicity of $V$ is assumed; the game may be arbitrarily non-monotone, subject only to the
quantitative bound implied by Lipschitz continuity.

\begin{lemma}[Bounded non-monotonicity]\label{lem:nonmono}
Under Assumption~\ref{as:lip}, for all $\boldsymbol{\theta}, \boldsymbol{\theta}' \in \Theta$,
\begin{equation}
\langle V(\boldsymbol{\theta}) - V(\boldsymbol{\theta}'), \boldsymbol{\theta} - \boldsymbol{\theta}' \rangle
\ge -L_V \|\boldsymbol{\theta} - \boldsymbol{\theta}'\|^2 .
\end{equation}
\end{lemma}

\begin{proof}
By Cauchy--Schwarz and Assumption~\ref{as:lip},
$\langle V(\boldsymbol{\theta}) - V(\boldsymbol{\theta}'), \boldsymbol{\theta} - \boldsymbol{\theta}' \rangle
\ge - \|V(\boldsymbol{\theta}) - V(\boldsymbol{\theta}')\| \, \|\boldsymbol{\theta} - \boldsymbol{\theta}'\|
\ge - L_V \|\boldsymbol{\theta} - \boldsymbol{\theta}'\|^2$.
\end{proof}

\begin{assumption}[Sufficient regularization]\label{as:lambda}
$\lambda > L_V / \rho_R$, where $\rho_R > 0$ is the constant computed in Lemma~\ref{lem:kl}.
\end{assumption}

\subsection*{3. Properties of the Gaussian KL Regularizer}

\begin{lemma}[Explicit regularizer constants]\label{lem:kl}
On $\Theta$ the operator $R$ is continuously differentiable, and
\begin{equation}
R(\boldsymbol{\theta}) = \begin{pmatrix}
\dfrac{\mu_x - \bar\mu_x}{\bar\sigma_x^{2}} \\[2ex]
-\dfrac{1}{\sigma_x} + \dfrac{\sigma_x}{\bar\sigma_x^{2}} \\[2ex]
\dfrac{\mu_y - \bar\mu_y}{\bar\sigma_y^{2}} \\[2ex]
-\dfrac{1}{\sigma_y} + \dfrac{\sigma_y}{\bar\sigma_y^{2}}
\end{pmatrix},
\qquad
J_R(\boldsymbol{\theta}) = \mathrm{diag}\!\left(
\frac{1}{\bar\sigma_x^{2}},\;
\frac{1}{\sigma_x^{2}} + \frac{1}{\bar\sigma_x^{2}},\;
\frac{1}{\bar\sigma_y^{2}},\;
\frac{1}{\sigma_y^{2}} + \frac{1}{\bar\sigma_y^{2}}
\right).
\end{equation}
Consequently $R$ is $\rho_R$-strongly monotone and $L_R$-Lipschitz on $\Theta$ with
\begin{equation}
\rho_R = \min\!\left(\frac{1}{\bar\sigma_x^{2}}, \frac{1}{\bar\sigma_y^{2}}\right) > 0,
\qquad
L_R = \frac{1}{\sigma_{\min}^{2}} + \max\!\left(\frac{1}{\bar\sigma_x^{2}}, \frac{1}{\bar\sigma_y^{2}}\right) < \infty .
\end{equation}
\end{lemma}

\begin{proof}
For univariate Gaussians,
\begin{equation}
D_{\mathrm{KL}}\big(\mathcal{N}(\mu,\sigma^2) \parallel \mathcal{N}(\bar\mu,\bar\sigma^2)\big)
= \log\frac{\bar\sigma}{\sigma} + \frac{\sigma^{2} + (\mu - \bar\mu)^{2}}{2 \bar\sigma^{2}} - \frac{1}{2},
\end{equation}
whose partial derivatives are $(\mu - \bar\mu)/\bar\sigma^{2}$ and $-1/\sigma + \sigma/\bar\sigma^{2}$, and whose
Hessian is $\mathrm{diag}(1/\bar\sigma^{2},\; 1/\sigma^{2} + 1/\bar\sigma^{2})$. Stacking the two players'
blocks gives the stated $R$ and $J_R$; the Jacobian is diagonal because the two KL terms are separable and
because $\bar{\boldsymbol{\theta}}$ is constant.

All four diagonal entries are positive and finite on $\Theta$, since $\sigma_x, \sigma_y \ge \sigma_{\min} > 0$.
As $1/\sigma^{2} + 1/\bar\sigma^{2} > 1/\bar\sigma^{2}$, the smallest entry is
$\min(\bar\sigma_x^{-2}, \bar\sigma_y^{-2}) = \rho_R$ and the largest is at most
$\sigma_{\min}^{-2} + \max(\bar\sigma_x^{-2}, \bar\sigma_y^{-2}) = L_R$. Hence
$\rho_R I \preceq J_R(\boldsymbol{\theta}) \preceq L_R I$ for every $\boldsymbol{\theta} \in \Theta$.

Since $\Theta$ is convex, the segment $\boldsymbol{\theta}' + t(\boldsymbol{\theta} - \boldsymbol{\theta}')$,
$t \in [0,1]$, lies in $\Theta$, and the fundamental theorem of calculus gives
\begin{equation}
\langle R(\boldsymbol{\theta}) - R(\boldsymbol{\theta}'), \boldsymbol{\theta} - \boldsymbol{\theta}' \rangle
= \int_0^1 (\boldsymbol{\theta} - \boldsymbol{\theta}')^{T}
J_R\big(\boldsymbol{\theta}' + t(\boldsymbol{\theta} - \boldsymbol{\theta}')\big)
(\boldsymbol{\theta} - \boldsymbol{\theta}') \, dt
\ge \rho_R \|\boldsymbol{\theta} - \boldsymbol{\theta}'\|^{2},
\end{equation}
and likewise
$\|R(\boldsymbol{\theta}) - R(\boldsymbol{\theta}')\|
= \big\| \int_0^1 J_R(\cdot)(\boldsymbol{\theta} - \boldsymbol{\theta}') dt \big\|
\le L_R \|\boldsymbol{\theta} - \boldsymbol{\theta}'\|$.
\end{proof}

\begin{lemma}[Strong monotonicity and Lipschitz continuity of $F$]\label{lem:F}
Under Assumptions~\ref{as:lip}--\ref{as:lambda}, $F = V + \lambda R$ is $\alpha$-strongly monotone and
$L_F$-Lipschitz on $\Theta$, with
\begin{equation}
\alpha = \lambda \rho_R - L_V > 0,
\qquad
L_F = L_V + \lambda L_R < \infty,
\qquad
\alpha \le L_F .
\end{equation}
\end{lemma}

\begin{proof}
Adding Lemma~\ref{lem:nonmono} to $\lambda$ times the strong monotonicity of Lemma~\ref{lem:kl} gives
\begin{equation}
\langle F(\boldsymbol{\theta}) - F(\boldsymbol{\theta}'), \boldsymbol{\theta} - \boldsymbol{\theta}' \rangle
\ge (\lambda \rho_R - L_V) \|\boldsymbol{\theta} - \boldsymbol{\theta}'\|^{2}
= \alpha \|\boldsymbol{\theta} - \boldsymbol{\theta}'\|^{2},
\end{equation}
and $\alpha > 0$ is exactly Assumption~\ref{as:lambda}. The triangle inequality applied to the Lipschitz
bounds of $V$ and $\lambda R$ gives $L_F = L_V + \lambda L_R$. Finally, combining the two properties with
Cauchy--Schwarz,
$\alpha \|\boldsymbol{\theta} - \boldsymbol{\theta}'\|^{2}
\le \langle F(\boldsymbol{\theta}) - F(\boldsymbol{\theta}'), \boldsymbol{\theta} - \boldsymbol{\theta}' \rangle
\le L_F \|\boldsymbol{\theta} - \boldsymbol{\theta}'\|^{2}$, so $\alpha \le L_F$.
\end{proof}

\subsection*{4. The Equilibrium: a Variational Inequality, not a Zero of $F$}

Because $\Theta$ has a boundary, the natural solution concept for \eqref{eq:pgda} is not $F(\boldsymbol{\theta}^*) = \mathbf{0}$
but the variational inequality $\mathrm{VI}(F, \Theta)$: find $\boldsymbol{\theta}^* \in \Theta$ with
\begin{equation}\label{eq:vi}
\langle F(\boldsymbol{\theta}^*), \boldsymbol{\theta} - \boldsymbol{\theta}^* \rangle \ge 0
\qquad \forall \boldsymbol{\theta} \in \Theta .
\end{equation}
Assuming a zero of $F$ inside $\Theta$ would be a genuine restriction: if the regularized equilibrium sits on
the boundary $\{\sigma_x = \sigma_{\min}\}$ or $\{\sigma_y = \sigma_{\min}\}$ --- precisely the variance-collapse
regime of interest --- no such zero exists, and a theorem premised on one would be vacuous.

\begin{proposition}[Existence, uniqueness, fixed-point characterization]\label{prop:vi}
Under Assumptions~\ref{as:lip}--\ref{as:lambda}, $\mathrm{VI}(F,\Theta)$ has exactly one solution
$\boldsymbol{\theta}^*$, and for every $\eta > 0$
\begin{equation}\label{eq:fixed}
\boldsymbol{\theta}^* = \Pi_\Theta\big(\boldsymbol{\theta}^* - \eta F(\boldsymbol{\theta}^*)\big).
\end{equation}
If moreover $\sigma_x^* > \sigma_{\min}$ and $\sigma_y^* > \sigma_{\min}$, then $F(\boldsymbol{\theta}^*) = \mathbf{0}$.
\end{proposition}

\begin{proof}
$\Theta$ is nonempty, closed and convex, and by Lemma~\ref{lem:F} $F$ is continuous and $\alpha$-strongly
monotone with $\alpha > 0$; strong monotonicity implies coercivity, so a solution exists and is unique
(standard theory of variational inequalities). For uniqueness directly: if $\boldsymbol{\theta}_1^*, \boldsymbol{\theta}_2^*$
both solve \eqref{eq:vi}, testing each with the other and adding gives
$\langle F(\boldsymbol{\theta}_1^*) - F(\boldsymbol{\theta}_2^*), \boldsymbol{\theta}_1^* - \boldsymbol{\theta}_2^* \rangle \le 0$,
while strong monotonicity forces this to be $\ge \alpha \|\boldsymbol{\theta}_1^* - \boldsymbol{\theta}_2^*\|^{2}$;
hence $\boldsymbol{\theta}_1^* = \boldsymbol{\theta}_2^*$.

The characterization \eqref{eq:fixed} is the standard projection identity: for a closed convex $\Theta$,
$z = \Pi_\Theta(w)$ iff $\langle w - z, \boldsymbol{\theta} - z\rangle \le 0$ for all $\boldsymbol{\theta} \in \Theta$.
Applying this with $w = \boldsymbol{\theta}^* - \eta F(\boldsymbol{\theta}^*)$ and $z = \boldsymbol{\theta}^*$ turns
\eqref{eq:fixed} into $\langle -\eta F(\boldsymbol{\theta}^*), \boldsymbol{\theta} - \boldsymbol{\theta}^*\rangle \le 0$,
which for $\eta > 0$ is exactly \eqref{eq:vi}.

If $\boldsymbol{\theta}^*$ lies in the interior of $\Theta$, then for small $\epsilon > 0$ every
$\boldsymbol{\theta} = \boldsymbol{\theta}^* \pm \epsilon e_i$ is feasible; \eqref{eq:vi} then gives both
$\pm \langle F(\boldsymbol{\theta}^*), e_i \rangle \ge 0$, so $F(\boldsymbol{\theta}^*) = \mathbf{0}$.
\end{proof}

\subsection*{5. Linear Convergence}

\begin{theorem}\label{thm:main}
Let Assumptions~\ref{as:lip}--\ref{as:lambda} hold, let $\boldsymbol{\theta}^*$ be the unique solution of
$\mathrm{VI}(F,\Theta)$ from Proposition~\ref{prop:vi}, and let $\alpha, L_F$ be as in Lemma~\ref{lem:F}. If
\begin{equation}
0 < \eta < \frac{2\alpha}{L_F^{2}},
\end{equation}
then the PGDA sequence \eqref{eq:pgda} converges to $\boldsymbol{\theta}^*$ linearly from any
$\boldsymbol{\theta}_0 \in \Theta$:
\begin{equation}
\|\boldsymbol{\theta}_k - \boldsymbol{\theta}^*\|^{2}
\le r(\eta)^{k} \|\boldsymbol{\theta}_0 - \boldsymbol{\theta}^*\|^{2},
\qquad
r(\eta) = 1 - 2\eta\alpha + \eta^{2} L_F^{2} \in [0, 1).
\end{equation}
\end{theorem}

\begin{proof}
Write $\boldsymbol{\delta}_k = \boldsymbol{\theta}_k - \boldsymbol{\theta}^*$.

\paragraph{Step 1: both iterates as projections.}
By \eqref{eq:pgda} and the fixed-point characterization \eqref{eq:fixed} --- used with the \emph{same} $\eta$ ---
\begin{equation}
\|\boldsymbol{\theta}_{k+1} - \boldsymbol{\theta}^*\|^{2}
= \big\| \Pi_\Theta\big(\boldsymbol{\theta}_k - \eta F(\boldsymbol{\theta}_k)\big)
- \Pi_\Theta\big(\boldsymbol{\theta}^* - \eta F(\boldsymbol{\theta}^*)\big) \big\|^{2}.
\end{equation}
This step is where the earlier reliance on $F(\boldsymbol{\theta}^*) = \mathbf{0}$ is removed: no interiority of
$\boldsymbol{\theta}^*$ is needed.

\paragraph{Step 2: non-expansiveness of the projection.}
$\Theta$ is closed and convex, so $\Pi_\Theta$ is non-expansive, giving
\begin{equation}
\|\boldsymbol{\theta}_{k+1} - \boldsymbol{\theta}^*\|^{2}
\le \big\| \boldsymbol{\delta}_k - \eta \big(F(\boldsymbol{\theta}_k) - F(\boldsymbol{\theta}^*)\big) \big\|^{2}.
\end{equation}

\paragraph{Step 3: expansion.}
\begin{equation}
\big\| \boldsymbol{\delta}_k - \eta \big(F(\boldsymbol{\theta}_k) - F(\boldsymbol{\theta}^*)\big) \big\|^{2}
= \|\boldsymbol{\delta}_k\|^{2}
- 2\eta \big\langle \boldsymbol{\delta}_k, F(\boldsymbol{\theta}_k) - F(\boldsymbol{\theta}^*) \big\rangle
+ \eta^{2} \big\| F(\boldsymbol{\theta}_k) - F(\boldsymbol{\theta}^*) \big\|^{2}.
\end{equation}

\paragraph{Step 4: bounding the two terms.}
By Lemma~\ref{lem:F} and $\eta > 0$,
\begin{align}
- 2\eta \big\langle \boldsymbol{\delta}_k, F(\boldsymbol{\theta}_k) - F(\boldsymbol{\theta}^*) \big\rangle
&\le -2\eta\alpha \|\boldsymbol{\delta}_k\|^{2}
&&\text{($\alpha$-strong monotonicity)}, \\
\eta^{2} \big\| F(\boldsymbol{\theta}_k) - F(\boldsymbol{\theta}^*) \big\|^{2}
&\le \eta^{2} L_F^{2} \|\boldsymbol{\delta}_k\|^{2}
&&\text{($L_F$-Lipschitz continuity)} .
\end{align}
Hence $\|\boldsymbol{\delta}_{k+1}\|^{2} \le r(\eta) \|\boldsymbol{\delta}_k\|^{2}$ with
$r(\eta) = 1 - 2\eta\alpha + \eta^{2} L_F^{2}$.

\paragraph{Step 5: the contraction factor.}
Since $\alpha \le L_F$ (Lemma~\ref{lem:F}),
\begin{equation}
r(\eta) = (1 - \eta \alpha)^{2} + \eta^{2}\big(L_F^{2} - \alpha^{2}\big) \ge 0,
\end{equation}
so $r(\eta) \in [0,1)$ as soon as $r(\eta) < 1$, i.e.\ $\eta(2\alpha - \eta L_F^{2}) > 0$, i.e.\
$0 < \eta < 2\alpha / L_F^{2}$. Iterating the one-step bound over $k$ steps gives the claim; $r(\eta) < 1$
also yields $\boldsymbol{\theta}_k \to \boldsymbol{\theta}^*$.
\end{proof}

\begin{corollary}[Optimal step size]
Minimizing $r(\eta)$ over $\eta$ gives $\eta^\star = \alpha / L_F^{2}$, which lies in $(0, 2\alpha/L_F^{2})$ and
achieves
\begin{equation}
r^\star = 1 - \frac{2\alpha^{2}}{L_F^{2}} + \frac{\alpha^{2}}{L_F^{2}}
= 1 - \frac{\alpha^{2}}{L_F^{2}} \in [0, 1).
\end{equation}
\end{corollary}

\begin{corollary}[Local version of Assumption~\ref{as:lip}]\label{cor:local}
Theorem~\ref{thm:main} remains valid if Assumption~\ref{as:lip} holds only on the bounded convex set
$B = \{\boldsymbol{\theta} \in \Theta : \|\boldsymbol{\theta} - \boldsymbol{\theta}^*\| \le \|\boldsymbol{\theta}_0 - \boldsymbol{\theta}^*\|\}$.
\end{corollary}

\begin{proof}
Induction on $k$. If $\boldsymbol{\theta}_k \in B$, then $\boldsymbol{\theta}_k, \boldsymbol{\theta}^* \in B$ and,
$B$ being convex, the segment joining them lies in $B$; all constants used in Steps 3--4 are therefore valid, so
$\|\boldsymbol{\delta}_{k+1}\| \le \sqrt{r(\eta)} \|\boldsymbol{\delta}_k\| \le \|\boldsymbol{\delta}_k\|$ and
$\boldsymbol{\theta}_{k+1} \in B$. The argument is not circular: each step only uses the constants on $B$, whose
definition does not depend on the iterates beyond $\boldsymbol{\theta}_0$.
\end{proof}

\subsection*{6. Scope of the Result}

\begin{remark}[The limit point is magnet-biased, and is not a Nash equilibrium]
$\boldsymbol{\theta}^*$ depends on both $\lambda$ and $\bar{\boldsymbol{\theta}}$: it solves
$\mathrm{VI}(V + \lambda R, \Theta)$, not $\mathrm{VI}(V, \Theta)$. The theorem therefore establishes convergence
of a single inner loop at a fixed magnet. Concluding anything about the unregularized game requires an outer
loop (e.g.\ periodically setting $\bar{\boldsymbol{\theta}} \leftarrow \boldsymbol{\theta}_k$, or annealing
$\lambda \downarrow 0$), whose analysis is not contained here.
\end{remark}

\begin{remark}[Restriction to the Gaussian family]
$\boldsymbol{\theta}^*$ is a stationary point of the regularized game \emph{within the parametric family of
Gaussian policies}. It coincides with the quantal-response equilibrium of the underlying continuous game only
when that equilibrium happens to be Gaussian; in general the two differ, and the result should not be stated as
convergence to ``the QRE''.
\end{remark}

\begin{remark}[Exact operator]
Iteration \eqref{eq:pgda} uses the exact $F(\boldsymbol{\theta}_k)$. Implementations estimate
$\nabla U$ from samples, giving $\boldsymbol{\theta}_{k+1} = \Pi_\Theta(\boldsymbol{\theta}_k - \eta \widehat{F}_k)$
with $\mathbb{E}[\widehat{F}_k] = F(\boldsymbol{\theta}_k)$ and variance $\varsigma^{2}$. The recursion then becomes
$\mathbb{E}\|\boldsymbol{\delta}_{k+1}\|^{2} \le r(\eta) \mathbb{E}\|\boldsymbol{\delta}_k\|^{2} + \eta^{2}\varsigma^{2}$,
so a constant step size yields linear convergence only to a noise ball of radius
$O(\eta \varsigma^{2} / \alpha)$; exact convergence requires decaying steps and a separate argument.
\end{remark}

\begin{remark}[Global Lipschitz continuity of $V$ is strong]
$\Theta$ is unbounded in $\mu_x, \mu_y$ and in the upper range of $\sigma_x, \sigma_y$. Assumption~\ref{as:lip}
holds globally for, e.g., bilinear $f(x,y) = xy$ (where $U = \mu_x \mu_y$ and $V = (\mu_y, 0, -\mu_x, 0)^{T}$),
but fails for generic polynomial payoffs. Corollary~\ref{cor:local} is the practical form of the assumption:
$L_V$ need only bound $V$ on the bounded set $B$ determined by the initialization.
\end{remark}

\begin{remark}[Parametrization]
The constants in Lemma~\ref{lem:kl} are for the direct parametrization by $\sigma$ with a hard projection onto
$\sigma \ge \sigma_{\min}$. An implementation parametrizing $s = \log \sigma$ (thereby making $\Theta$ all of
$\mathbb{R}^4$ and dispensing with the projection) induces a different Jacobian, and $\rho_R$, $L_R$ must be
recomputed accordingly.
\end{remark}

\end{document}
